% Cote 119: offprints, reports, courses, CV and letters (typed and printed part): archive cote 119, site page 197 (PDF page 198).
% typescript. Draft compared: none.
% AI-transcribed by Claude Opus 5.5 (claude-opus-5-5, adaptive thinking, effort high), Message Batches API, batch r119_full (msgbatch_01DwZNm5Y7YsKBWyAN4sFs8w), 2026-11-13; not yet checked by eye.
\scanpage{119/197}
\mspage{- 2 -}

2. \uline{Algèbre Homologique} ([8], [9], [19], [9 bis])

Depuis 1955, me plaçant au point de vue de ``l'usager'' et non celui
du spécialiste, j'ai été amené continuellement à élargir et à assouplir le
langage de l'algèbre homologique, notamment sous la poussée des besoins de la
géométrie algébrique (théories de dualité, théories du type Riemann-Roch,
cohomologies $\ell$-adiques, cohomologies du type de De Rham, cohomologies cristal-%
lines ...). Deux directions principales à ces réflexions  : développement d'une
algèbre homologique non commutative (amorcée dans [10 bis] et systématisée
dans la thèse de J. Giraud [*]); théorie des catégories dérivées (développée
systématiquement par J.L. Verdier, exposée dans Hartshorne [*], Illusie [*]
et [16 SGA 4 Exp. XVIII]). Ces deux courants de réflexion sont d'ailleurs loin
d'être épuisés, et sont sans doute appelés à se rejoindre,soit au sein d'une
``algèbre homotopique'' dont une esquisse préliminaire a été faite par Quillen
[*], soit dans l'esprit de la théorie des n-catégories , particulièrement bien
adaptée à l'interprétation géométrique des invariants cohomologiques (cf. le
livre cité de J. Giraud et le travail de Mme. M. Raynaud [*]).

\uline{Références} : J.L. Verdier, P. Deligne, D. Quillen, P. Gabriel.

3. \uline{Topologie} ([16, SGA 4], [9])

Jusqu'à présent, c'est surtout le K-invariant des espaces topologiques
que j'avais introduit à l'occasion de mes recherches sur le théorème de Riemann-%
Roch en géométrie algébrique, qui a connu la fortune\unsure{,} la plus brillante, étant
le point de départ de très nombreuses recherches en topologie homotopique et
topologie différentielle. De nombreuses constructions que j'avais introduites
pour les besoins de la démonstration algébrique du théorème de Riemann-Roch
(telles les opérations  $\lambda_i$  et leurs liens avec les opérations du groupe symé-%
trique) sont devenues pratique courante non seulement en géométrie algébrique
et en algèbre, mais également en topologie et en théorie des nombres, notamment
dans les travaux de mathématiciens comme Atiyah, Hirzebruch, Adams, Quillen,
Bass, Tate, Milnor, Karoubi, Shih,etc..

Plus fondamental me semble néanmoins  l'élargissement de la
topologie générale, dans l'esprit de la théorie des faisceaux (développée
initialement par J. Leray), contenu dans le point de vue des topos
[16, SGA 4]. J'ai introduit ces topos à partir de 1958 en partant du
besoin de définir une cohomologie  $\ell$-adique  des variétés algébriques

\mspage{\hand{197}}
